\documentclass[11pt]{article}
\usepackage[utf8]{inputenc}
\usepackage[margin=1in]{geometry}
\usepackage{amsmath}
\usepackage{graphicx}
\usepackage{booktabs}
\usepackage{hyperref}
\usepackage[round]{natbib}

\title{Emotional Vocalizations Alter Behaviors and Neurochemical Release into the Amygdala}
\author{Author A,$^{1}$ Author B,$^{1}$ Author C,$^{2}$ Author D$^{1,*}$\\
\small $^{1}$Department of Neurobiology, Northeast Ohio Medical University\\
\small $^{2}$Department of Anatomy, Northeast Ohio Medical University\\
\small $^{*}$Corresponding author: authord@neomed.edu}
\date{}

\begin{document}
\maketitle

\begin{abstract}
Social vocalizations convey emotional information that shapes behavioral decisions, yet the neurochemical mechanisms by which the brain differentially processes appetitive versus aversive vocal signals remain poorly understood. Here we used in vivo microdialysis in the basolateral amygdala (BLA) of CBA/CaJ mice combined with playback of mating and restraint vocalizations to measure acetylcholine (ACh) and dopamine (DA) release during emotional vocal processing. In experienced males, restraint vocalizations increased ACh and decreased DA release, while mating vocalizations produced the opposite pattern---decreased ACh and increased DA. Non-estrus females responded to mating vocalizations with male-like neurochemical patterns, while estrus females showed a unique profile with both ACh and DA increased. Critically, every group exhibiting increased ACh release also displayed increased aversive flinching behavior, suggesting a mechanistic link between cholinergic signaling and defensive responses. Serotonin metabolite (5-HIAA) showed no context-dependent changes, confirming the specificity of ACh/DA modulation. Inexperienced mice lacking prior behavioral experience with mating and restraint showed no distinct neurochemical patterns, demonstrating that a single 90-minute experience with each context is necessary and sufficient to establish these responses. These findings reveal that ACh and DA provide opposing, experience-dependent contextual signals to the BLA that modulate behavioral responses to emotional vocalizations.
\end{abstract}

\section{Introduction}

Social communication through vocalizations is fundamental to survival and reproduction across mammalian species \citep{portfors2007types}. Mice produce a rich repertoire of vocalizations in different behavioral contexts: ultrasonic vocalizations (USVs) during mating interactions \citep{holy2005ultrasonic}, and mid-frequency vocalizations (MFVs) and distress calls during aversive experiences such as restraint. These vocalizations carry emotional valence information that listeners must decode to generate appropriate behavioral responses---approaching a potential mate versus avoiding a threat.

The basolateral amygdala (BLA) is a critical node for integrating sensory information with emotional significance \citep{ledoux2000emotion, janak2015amygdala}. The BLA receives auditory input and combines it with contextual and internal-state information to shape behavioral output through projections to the central amygdala (CeA) for defensive responses and to the nucleus accumbens (NAc) for appetitive responses \citep{beyeler2016divergent, tye2013amygdala}. Neuromodulatory inputs, particularly acetylcholine from the basal forebrain \citep{ballinger2016basal, hasselmo2006role} and dopamine from the ventral tegmental area \citep{schultz2007multiple}, are well positioned to modulate BLA processing of emotional information \citep{mcgaugh2004amygdala}.

While the roles of ACh and DA in vocalization production have been studied, their contribution to the processing and interpretation of social vocalizations at the level of the BLA remains unexplored. Furthermore, internal state variables---sex, hormonal status, and prior experience---are known to influence amygdala processing, yet how these factors modulate the neurochemical response to vocal signals has not been investigated. Finally, whether a single behavioral experience is sufficient to establish context-dependent neurochemical responses, or whether repeated exposure is required, is unknown.

Here we address these questions using in vivo microdialysis to measure ACh, DA, and 5-HIAA release in the BLA during playback of mating and restraint vocal sequences in experienced and inexperienced male and female CBA/CaJ mice.

\section{Methods}

\subsection{Animals}

CBA/CaJ mice ($n = 83$; p90--p180; both sexes; Jackson Labs) were housed in same-sex groups on a reversed dark/light cycle. Experiments were conducted during the dark phase. Female estrous stage was determined by vaginal smear cytology (crystal violet staining) before and after experiments.

\subsection{Experimental Timeline}

The 6-day protocol was as follows: Days 1--2, each mouse underwent one bout each of mating behavior (high-intensity interactions including mounting) and sustained restraint (90~minutes each, counterbalanced order). After the Day 2 experience, a microdialysis guide tube was implanted above the BLA \citep{paxinos2001mouse}. Days 3--5, recovery. Day 6, microdialysis probe insertion, equilibration, and vocal playback with sampling.

\begin{table}[ht]
\centering
\caption{Experimental Groups}
\label{tab:groups}
\begin{tabular}{llllc}
\toprule
\textbf{Group} & \textbf{Sex} & \textbf{Experience} & \textbf{Playback} & \textbf{$n$} \\
\midrule
EXP Male -- Mating & Male & Experienced & Mating & 7 \\
EXP Male -- Restraint & Male & Experienced & Restraint & 6 \\
EXP Estrus Female & Female (estrus) & Experienced & Mating & 6 \\
EXP Non-estrus Female & Female (non-estrus) & Experienced & Mating & 5 \\
INEXP Male -- Mating & Male & Inexperienced & Mating & 7 \\
INEXP Male -- Restraint & Male & Inexperienced & Restraint & 7 \\
INEXP Estrus Female & Female (estrus) & Inexperienced & Mating & 7 \\
\bottomrule
\end{tabular}
\end{table}

\subsection{Vocal Stimuli}

Mating vocal sequences contained 545 syllables (primarily USVs with harmonics, steps, and complex structure from males; LFH calls from females). Restraint vocal sequences contained 622 syllables (primarily MFVs from restrained mice). Each 20-minute playback featured 7 repetitions of selected sequences from a single vocalization type.

\subsection{Microdialysis and Neurochemical Analysis}

Microdialysis samples were collected in 10-minute intervals (Pre-Stim, Stim 1, Stim 2) and analyzed by liquid chromatography/mass spectrometry (LC/MS) for ACh, DA, 5-HIAA, and other neurochemicals. Each animal participated in a single playback session. Probe placement was verified post-experiment by infusion of fluorescent dextran tracers followed by histological confirmation.

\subsection{Behavioral Scoring}

Behavioral responses were manually scored from video recordings during each 10-minute sampling period. Behaviors included attending (orienting toward sound source), still-and-alert (freezing with apparent vigilance), flinching (startle-like reflexive movement), grooming, rearing, and walking.

\section{Results}

\subsection{Opposing ACh and DA Patterns in Males}

Experienced male mice showed striking, opposing neurochemical responses to mating versus restraint vocalizations (Table~\ref{tab:male_neuro}).

\begin{table}[ht]
\centering
\caption{Neurochemical Responses in Experienced Males}
\label{tab:male_neuro}
\begin{tabular}{lccc}
\toprule
\textbf{Neurochemical} & \textbf{Restraint Playback} & \textbf{Mating Playback} & \textbf{5-HIAA} \\
\midrule
ACh & Increased & Decreased & No change \\
DA & Decreased & Increased & No change \\
\bottomrule
\end{tabular}
\end{table}

Restraint vocalizations increased ACh and decreased DA release in the BLA during both Stim 1 and Stim 2 periods. Conversely, mating vocalizations decreased ACh and increased DA release. 5-HIAA (serotonin metabolite) showed no context-dependent changes in either condition, confirming the specificity of the ACh/DA findings.

Behaviorally, restraint vocalizations increased still-and-alert and flinching behaviors in males, while mating vocalizations produced no significant increases in aversive behaviors. Both vocalization types increased attending behavior during Stim 1.

\subsection{Sex and Hormonal State Modulate Neurochemical Responses}

Analysis of mating vocal playback responses across sex and hormonal state revealed a critical dissociation (Table~\ref{tab:sex}).

\begin{table}[ht]
\centering
\caption{ACh and DA Responses to Mating Playback by Sex and Estrous State}
\label{tab:sex}
\begin{tabular}{lcc}
\toprule
\textbf{Group} & \textbf{ACh} & \textbf{DA} \\
\midrule
Male ($n = 7$) & Decreased & Increased \\
Non-estrus female ($n = 5$) & Decreased (male-like) & Increased (male-like) \\
Estrus female ($n = 6$) & Increased (unique) & Increased \\
\bottomrule
\end{tabular}
\end{table}

Non-estrus females showed neurochemical patterns indistinguishable from males in response to mating vocalizations: decreased ACh and increased DA. Estrus females, however, showed a unique profile with increased ACh alongside increased DA. This estrus-specific ACh increase represents a state-dependent modulation of vocal processing by the reproductive cycle.

Attending behavior increased during Stim 1 regardless of sex or estrous stage (repeated-measures ANOVA: time effect $F(2,30) = 32.6$, $p < 0.001$, $\eta^2 = 0.7$; time $\times$ sex interaction $F(2,30) = 0.12$, $p = 0.9$, $\eta^2 = 0.008$).

\subsection{ACh Release Co-Occurs with Aversive Flinching Behavior}

A consistent pattern emerged across all experimental groups: every group showing significantly increased ACh release also displayed significantly increased flinching behavior, and this co-occurrence was present during both Stim 1 and Stim 2 periods (Table~\ref{tab:flinch}).

\begin{table}[ht]
\centering
\caption{Co-Occurrence of ACh Release and Flinching Across Groups}
\label{tab:flinch}
\begin{tabular}{lcc}
\toprule
\textbf{Group} & \textbf{ACh Change} & \textbf{Flinching Change} \\
\midrule
Male -- Restraint & Increased & Increased \\
Male -- Mating & Decreased & No change \\
Estrus Female -- Mating & Increased & Increased \\
Non-estrus Female -- Mating & Decreased & No change \\
\bottomrule
\end{tabular}
\end{table}

This temporal matching across multiple independent groups with different stimuli and internal states provides compelling correlational evidence for a mechanistic relationship between cholinergic BLA input and aversive behavioral responses.

\subsection{A Single Experience Is Necessary and Sufficient}

Inexperienced (INEXP) mice---those that underwent the same protocol but without prior mating or restraint experience---showed no distinct context-dependent neurochemical patterns. ACh, DA, and 5-HIAA all failed to show significant changes from baseline in response to either vocalization type. Furthermore, flinching and still-and-alert behaviors were not increased in INEXP mice, although attending behavior still increased (indicating intact auditory processing).

There were no pre-stimulus baseline differences between EXP and INEXP mice for any behavior or neurochemical, confirming that the experience manipulation specifically affected stimulus-evoked responses, not basal state. This demonstrates that a single 90-minute bout each of mating and restraint experience is necessary and sufficient to establish the context-dependent neurochemical response patterns.

\section{Discussion}

This study reveals that the BLA receives context-dependent, opposing neuromodulatory signals during the processing of emotional vocalizations: ACh increases during aversive (restraint) vocal processing while DA increases during appetitive (mating) vocal processing. These patterns are experience-dependent, sex- and hormonal-state-modulated, and co-occur with specific behavioral responses.

The opposing ACh/DA pattern is consistent with the known functional roles of these neuromodulators in the amygdala circuit. Cholinergic input from the basal forebrain \citep{ballinger2016basal} enhances excitatory responses in CeA-projecting BLA neurons through muscarinic M1 receptors, promoting defensive behavioral output. Dopaminergic input from the VTA \citep{schultz2007multiple} enhances responses in NAc-projecting BLA neurons through D1 receptors, promoting appetitive behavioral output \citep{beyeler2016divergent}. Our data suggest that these two neuromodulatory systems work in opposition to bias BLA output toward the appropriate behavioral response.

The sex- and hormonal-state-dependent modulation is particularly striking. The finding that non-estrus females respond to mating vocalizations with a male-like pattern (decreased ACh, increased DA) while estrus females show increased ACh suggests that the hormonal state of the listener actively reconfigures the neurochemical interpretation of conspecific vocal signals. The estrus-specific ACh increase may reflect heightened vigilance or arousal associated with the reproductive state.

The experience-dependence finding has important implications for understanding how vocal communication develops. A single 90-minute experience with both mating and restraint contexts was sufficient to establish robust, stimulus-specific neurochemical responses---but without this experience, vocal playback evoked no distinct patterns. This suggests that the BLA requires associative learning between vocal signals and their behavioral contexts before it can generate appropriate neurochemical responses.

Limitations include the use of microdialysis, which provides temporal resolution of minutes rather than seconds, and the inability to determine whether ACh and DA act on the same or different BLA neurons. Future studies using fiber photometry or cell-type-specific sensors could address these questions.

\section{Conclusion}

We demonstrate that emotional vocalizations elicit opposing, context-dependent patterns of ACh and DA release in the basolateral amygdala: aversive restraint vocalizations increase ACh and decrease DA, while appetitive mating vocalizations decrease ACh and increase DA. These patterns are modulated by sex and estrous state, co-occur with specific behavioral responses (ACh with flinching), and require prior behavioral experience to emerge. A single experience with mating and restraint contexts is both necessary and sufficient. These findings establish ACh and DA as complementary contextual signals that shape BLA output during emotional vocal communication.

\bibliographystyle{plainnat}
\bibliography{references}

\end{document}
